% =====================================================================
% CHAPTER 2 (9-chapter template, RESTRUCTURE_PLAN_V2): Literature Review
%   2.1 Introduction (rewritten roadmap for the new 4-section shape)
%   2.2 Existing Research        <- old ch2 2.2 Literature Review
%                                   + old ch2 2.3 Requirements of
%                                   Edge-Deployable Stereo Matching
%                                   + old ch2 2.4 Classification of
%                                   Efficient Stereo Methods, all folded
%                                   into flat \subsection entries (no
%                                   sub-subsections)
%   2.3 Comparative Analysis     <- old ch2 2.5 Cross-Domain
%                                   Generalization and Benchmarks
%                                   + the comparison table/figure lifted
%                                   out of old ch2 2.6, reframed as a
%                                   direct family-vs-family comparison
%   2.4 Summary                  <- old ch2 2.7, updated to the new
%                                   shape, handing off to Chapter 3 via
%                                   \ref{ch:design}
% Old ch2 2.6 "Detailed Review of the Adopted Methodology" (the four
% adopted-mechanism narrative) is EXCLUDED here; it is owned by the ch3
% agent. All \label names carried over from the old chapter are kept
% verbatim; equations are byte-identical; every hard-coded "Chapter N"
% has been replaced with a \ref to the corresponding new chapter label.
% =====================================================================
\chapter{Literature Review}
\label{ch:litreview}
\label{ch:lit}

\section{Introduction}
This chapter surveys the literature that shapes the design of this
thesis. Section~\ref{sec:existing_research} traces deep stereo matching
chronologically, from the classical four-step pipeline and the first
learned matching costs through three-dimensional cost-volume networks,
iterative refinement, and the foundation-model systems that currently
hold the accuracy ceiling. Because this thesis targets deployment on an
edge accelerator rather than accuracy alone, that section then states
the hardware requirements edge deployment imposes and classifies the
efficient and edge-deployable literature by the technique family each
uses to meet them. Section~\ref{sec:research_gap} draws the survey
together into the specific gap this thesis targets.
Section~\ref{sec:comparative_analysis} then compares the families
directly: it reviews cross-domain generalization and the benchmarks used
to measure it, then sets published parameter counts, accuracy, and
latency figures for the surveyed methods side by side, alongside our own
model. Section~\ref{sec:ch2_summary} closes the chapter with the design
lessons this survey establishes.

\FloatBarrier
\section{Existing Research}
\label{sec:existing_research}

\subsection{Stereo Matching Fundamentals}
Stereo matching is the problem of finding, for each pixel in one image of
a stereo pair, the pixel in the other image that captures the same scene
point \cite{scharstein2002taxonomy}. This subsection sets out the concepts
that the rest of the thesis relies on: the epipolar geometry that reduces
the search to one dimension, rectification, the disparity that encodes
depth, the matching costs that measure image agreement, the cost volume
that organizes those costs, the local and global families that solve for
disparity, and the refinement that cleans the result. These ideas were
formalized in the classical era \cite{scharstein2002taxonomy} and survive,
in learned form, in the deep networks reviewed later in this chapter
\cite{poggi2021synergies, laga2020survey}.

\paragraph{Epipolar geometry.}
The search for a match is constrained by geometry rather than left open
over the whole image. Given the relative pose of the two cameras, a point
seen at position $p$ in the left image must lie, in the right image, on a
single straight line called the epipolar line, and every epipolar line
passes through a fixed point called the epipole, the image of the other
camera's center. Figure \ref{fig:epipolar} illustrates this epipolar
constraint: it reduces the two-dimensional problem of searching the
entire right image to a one-dimensional problem of searching along a
known line.

\begin{figure}[H]
\centering
\includegraphics[width=0.70\textwidth]{fig_2_epipolar}
\caption{Epipolar geometry constraining stereo correspondence search to one dimension (Fei-Fei Li, CS231A).}
\label{fig:epipolar}
\end{figure}

\paragraph{Rectification.}
The epipolar lines of a general camera pair are slanted and vary from
row to row, which makes them difficult to search directly. Rectification
addresses this by reprojecting both images onto a common image plane parallel
to the line joining the two camera centers, as Figure
\ref{fig:rectify_geom} shows. After this warp every epipolar line
coincides with an image row, so corresponding points share the same
vertical coordinate and the match for a left pixel $(x, y)$ is found by
sliding horizontally \cite{scharstein2002taxonomy}:
\begin{equation}
x' = x - d, \qquad y' = y
\label{eq:correspondence}
\end{equation}
where $(x', y')$ is the matching pixel in the right image and $d \ge 0$ is
the disparity. We assume rectified input throughout, so the
correspondence search always reduces to a horizontal scan, and the
accuracy of the rectification itself becomes a source of error worth
examining directly.

\begin{figure}[H]
\centering
\includegraphics[width=0.30\textwidth]{fig_2_rectification}
\caption{Stereo rectification reprojecting tilted image planes onto a common plane (Fei-Fei Li, CS231A).}
\label{fig:rectify_geom}
\end{figure}

\paragraph{Disparity and depth.}
The horizontal shift $d$ between a point and its match is the disparity,
and it fixes the depth of the scene point through the calibrated pinhole
relation \cite{scharstein2002taxonomy}
\begin{equation}
Z = \frac{f\,B}{d}
\label{eq:depth_from_disparity}
\end{equation}

\begin{figure}[H]
\centering
\includegraphics[width=0.92\textwidth]{fig_1_1_stereo_geometry_preview}
\caption{Two-view stereo geometry relating disparity to depth.}
\label{fig:stereo_geometry}
\end{figure}

where $Z$ is the depth of a scene point, $f$ the camera focal length in
pixels, $B$ the baseline separating the two cameras, and $d$ the
disparity, the horizontal shift of the point between the two rectified
views (Figure \ref{fig:stereo_geometry}). Because the relation is inverse,
disparity resolution matters most for distant surfaces: a fixed sub-pixel
error in disparity corresponds to an increasingly large depth interval as
the scene recedes, and this is why the field reports accuracy in fractions
of a pixel.

\paragraph{The matching cost.}
To decide which candidate along the scanline is the true match, a stereo
method needs a numerical measure of how well two image locations agree,
called the matching cost. The classical costs trade robustness against
simplicity \cite{hirschmuller2007evalcosts}. The squared and absolute
intensity differences are cheap but assume the two cameras record the same
brightness, so they fail under gain and exposure differences. Normalized
cross-correlation over a window removes a constant gain and bias and is
therefore more stable. The census transform \cite{zabih1994census} is the
most robust of the classical costs: it replaces each pixel by a binary
string that records only whether each neighbor is brighter or darker than
the center, and it compares two such strings by their Hamming distance,
\begin{equation}
C_{\text{cen}}(p, d) = \mathrm{Ham}\!\big(\mathcal{T}(I_L, p),\ \mathcal{T}(I_R, p - d)\big),
\qquad
\mathcal{T}(I, p) = \bigotimes_{q \in N_p} \big[\, I(q) \ge I(p)\,\big]
\label{eq:census}
\end{equation}
where $\mathcal{T}$ is the census string built by concatenating (the
operator $\bigotimes$) the comparison bit $[\,I(q) \ge I(p)\,]$ over the
neighborhood $N_p$, $\mathrm{Ham}$ counts the positions at which two
strings differ, and $[\cdot]$ is the indicator function. Because it keeps
only the ordering of intensities and discards their magnitudes, the census
cost is insensitive to illumination changes between the two views, as
Figure \ref{fig:census} illustrates. This step is exactly what the first
deep methods replaced with a learned network: MC-CNN
\cite{zbontar2016mccnn} swapped the hand-crafted cost for a convolutional
similarity score while leaving the rest of the classical pipeline in
place, and the feature encoders of every modern network build on this
same idea.

\begin{figure}[H]
\centering
\includegraphics[width=0.82\textwidth]{fig_2_census}
\caption{The census transform for illumination-robust matching cost computation.}
\label{fig:census}
\end{figure}

\paragraph{The disparity space image and cost volume.}
Evaluating the matching cost for every pixel and every candidate disparity
produces a three-dimensional structure $C(x, y, d)$ called the disparity
space image, or cost volume: for each pixel it records how well the two
views agree at each possible shift, as visualized in Figure
\ref{fig:costvolume}. Reading the volume along the disparity axis at one
pixel gives a cost profile whose minimum is the winning disparity; slicing
it at a fixed disparity shows which regions of the scene sit at that depth.
This structure is the direct precursor of the learned cost volumes used
throughout this thesis: modern networks build the same volume from learned
features instead of raw intensities, but the object being reasoned over is
still $C(x, y, d)$.

\begin{figure}[H]
\centering
\includegraphics[width=0.60\textwidth]{fig_2_costvolume_times}
\caption{The disparity space image, or cost volume, for stereo matching.}
\label{fig:costvolume}
\end{figure}

\paragraph{The four-step pipeline.}
Scharstein and Szeliski \cite{scharstein2002taxonomy} showed that every
stereo algorithm, classical or deep, decomposes into the same four steps,
shown in Figure \ref{fig:classical_pipeline}: matching cost computation,
which fills the cost volume; cost aggregation, which smooths it by pooling
evidence from neighboring pixels; disparity optimization, which selects a
disparity at each pixel; and disparity refinement, which post-processes the
result. The rest of the chapter follows this decomposition, because deep
learning progressively replaced each of the four steps with a learned
module while keeping the overall structure intact.

\begin{figure}[H]
\centering
\includegraphics[width=0.82\textwidth]{fig_2_1_classical_preview}
\caption[The four-step classical stereo matching pipeline.]{The four-step classical stereo matching pipeline \cite{scharstein2002taxonomy}.}
\label{fig:classical_pipeline}
\end{figure}

\paragraph{Local and global methods.}
Classical algorithms split into two families by where they place their
effort. Local methods concentrate on aggregation: they average the
matching cost over a support window and then pick, at each pixel
independently, the disparity of lowest cost. This is fast and simple, but
it enforces smoothness only implicitly through the window, so it blurs
depth edges (the window straddles two depths at a boundary) and fails in
textureless regions where the window contains no distinctive signal.
Global methods instead concentrate on optimization: they define an energy
over the whole disparity map that adds a data term, rewarding low matching
cost, to a smoothness term, penalizing disparity jumps between neighbors,
and they search for the disparity map of least energy. This yields sharp,
consistent results because the smoothness prior is explicit and can be made
discontinuity-preserving, but the joint optimization is far more expensive
than per-pixel selection. The trade-off between the two families,
aggregation speed against optimization quality, reappears in the
efficient deep methods surveyed later in this chapter.

\paragraph{Disparity optimization and semi-global matching.}
Semi-global matching \cite{hirschmuller2008sgm} became the standard
classical algorithm by approximating the expensive two-dimensional global
optimization as a sum of independent one-dimensional dynamic-programming
problems solved along several path directions, which recovers most of the
quality of a full global method at a fraction of the cost. It minimizes an energy that
combines a per-pixel matching cost with two smoothness penalties over
neighboring pixels,
\begin{equation}
E(D) = \sum_{p} C(p, D_p)
  + \sum_{p} \sum_{q \in N_p} \Big[
      P_1\, \mathbb{1}\!\left[\,|D_p - D_q| = 1\,\right]
    + P_2\, \mathbb{1}\!\left[\,|D_p - D_q| > 1\,\right] \Big]
\label{eq:sgm_energy}
\end{equation}
where $D$ is the disparity map, $C(p, D_p)$ the matching cost of assigning
disparity $D_p$ to pixel $p$, $N_p$ the neighbors of $p$, $P_1$ a small
penalty for one-pixel disparity changes that preserves slanted surfaces,
$P_2$ a larger penalty for jumps beyond one pixel that permits true
discontinuities, and $\mathbb{1}[\cdot]$ the indicator function. Semi-global
matching links directly to deep stereo: the first learned method kept it
as an optimization backend, and later networks replaced its fixed
aggregation with learnable layers.

\paragraph{Disparity refinement.}
The disparity selected from a discrete cost volume is an integer, and it
still carries errors in occluded and textureless regions, so a final
refinement step improves both its precision and its reliability. Sub-pixel
precision is recovered by fitting a parabola \cite{scharstein2002taxonomy}
to the three cost values around the winning disparity $d$ and taking the
location of its minimum,
\begin{equation}
d_{\text{sub}} = d - \frac{C(d + 1) - C(d - 1)}{2\,\big[\,C(d + 1) - 2\,C(d) + C(d - 1)\,\big]}
\label{eq:subpixel}
\end{equation}
where $C(d)$ is the aggregated cost at the winning disparity and
$C(d \pm 1)$ its two neighbors, which yields a fractional disparity without
enlarging the search range. A left-right consistency check computes the
disparity from both images and flags as occluded any pixel whose two
estimates disagree by more than a pixel, and the flagged pixels are then
filled from nearby reliable values and cleaned by median filtering. Each of
these operations reappears in learned form in the modern networks reviewed
below, where refinement is carried out by trained convolutional or
recurrent modules rather than fixed rules.

\FloatBarrier
\subsection{The Deep Learning Transition}
Deep learning entered stereo matching first through the matching-cost
step. MC-CNN \cite{zbontar2016mccnn} replaced only the cost computation
with a Siamese convolutional network that scores the similarity of two
patches, and fed those scores into an otherwise classical pipeline of
cross-based aggregation, semi-global optimization, and refinement. It
established a trick that later networks kept: compute each image's
features once and reuse them across all disparities. The method itself
was not end-to-end. DispNetC \cite{mayer2016sceneflow} removed the
classical stages entirely and regressed disparity directly with an
encoder-decoder network. Its central component was a correlation layer,
which forms a compact cost volume by correlating the two feature maps
along the scanline, as in \eqref{eq:correlation}; it ran roughly a
thousand times faster than MC-CNN at comparable accuracy. The same work
introduced the SceneFlow synthetic dataset, which remains the standard
pretraining set for stereo networks. GC-Net \cite{kendall2017gcnet}
completed the transition to a fully learned cost-volume network. It
builds a four-dimensional cost volume by concatenating the left and right
feature vectors at each disparity, regularizes it with a
three-dimensional convolutional encoder-decoder, and regresses a
continuous disparity with the differentiable soft-argmin of
\eqref{eq:softargmin}, an operator still used across cost-volume networks
today.

The correlation layer \cite{mayer2016sceneflow} forms a scalar matching
cost by taking the inner product of the two feature maps at each
disparity along the scanline,
\begin{equation}
c(x, y, d) = \big\langle\, f_L(x, y),\; f_R(x - d,\, y) \,\big\rangle
           = \sum_{k} f_L^{(k)}(x, y)\, f_R^{(k)}(x - d, y)
\label{eq:correlation}
\end{equation}
where $c(x, y, d)$ is the matching cost for reference pixel $(x, y)$ at
disparity $d$, $f_L$ and $f_R$ the shared feature maps, and the inner
product runs over the feature channels $k$; the search covers
$d \in [0, D_{\max}]$.

The soft-argmin \cite{kendall2017gcnet} turns the regularized cost volume
into a continuous disparity by taking a softmax-weighted expectation,
\begin{equation}
\hat{d}(x, y) = \sum_{d=0}^{D-1} d \cdot
  \operatorname{softmax}_d\!\big(-c_d(x, y)\big)
\label{eq:softargmin}
\end{equation}
where $\hat{d}(x, y)$ is the regressed sub-pixel disparity and $c_d(x, y)$
the aggregated cost at disparity $d$; the negative sign turns low cost
into high probability. The operator assumes a single-peaked cost
distribution and returns an intermediate value where that assumption
fails, such as at a depth discontinuity.

\FloatBarrier
\subsection{The 3D Cost-Volume Era}
The years that followed refined the three-dimensional cost-volume
paradigm. PSMNet \cite{chang2018psmnet}, shown in Figure
\ref{fig:psmnet}, added spatial pyramid pooling to inject global context
at the feature level, along with three stacked hourglass modules under
intermediate supervision, and became the de-facto baseline for the era.

\begin{figure}[!htbp]
\centering
\includegraphics[width=\textwidth]{fig_2_4a_psmnet}
\caption[PSMNet architecture, representing the three-dimensional cost-volume paradigm.]{PSMNet architecture, representing the three-dimensional cost-volume paradigm \cite{chang2018psmnet}.}
\label{fig:psmnet}
\end{figure}

GwcNet \cite{guo2019gwcnet} introduced the group-wise correlation
volume, which splits the features into channel groups and computes an
inner-product correlation per group, as in \eqref{eq:gwc_lit}. This form
sits between the information-poor full correlation and the expensive
concatenation, and it degrades far more gracefully than concatenation
under a reduced aggregation budget, which is why nearly every later
efficient method adopts it, including the tile initialization we use.
GA-Net \cite{zhang2019ganet} addressed the cost of
three-dimensional convolution directly: it replaced the convolution with
a differentiable semi-global aggregation layer that approximates
recursive aggregation along a few directions and a local guided
aggregation layer for thin structures, and showed that a network with
only two such layers can outperform GC-Net with nineteen
three-dimensional convolutions. Later work pushed efficiency further:
AANet \cite{xu2020aanet} removed
three-dimensional convolution entirely in favor of deformable intra-scale
and cross-scale aggregation, ACVNet \cite{xu2022acvnet} used an
attention map derived from correlation to filter a concatenation volume,
and CFNet \cite{shen2021cfnet} shifted the objective toward cross-domain
generalization through a fused multi-scale volume and an
uncertainty-driven cascade that adaptively narrows the disparity search.

The group-wise correlation \cite{guo2019gwcnet} splits the $N_c$ feature
channels into $N_g$ groups and computes a normalized inner-product
correlation within each group,
\begin{equation}
C_{\mathrm{gwc}}(g, d, x, y) = \frac{N_g}{N_c}\,
  \big\langle\, f_L^{(g)}(x, y),\; f_R^{(g)}(x - d,\, y) \,\big\rangle
\label{eq:gwc_lit}
\end{equation}
where $g$ indexes the group, $d$ the disparity, and $f_L^{(g)}$,
$f_R^{(g)}$ the channel slice of each feature map belonging to group $g$;
the resulting volume carries one channel per group. We use this operator
at a single coarse scale in our tile initialization.

\FloatBarrier
\subsection{Iterative Refinement}
RAFT-Stereo \cite{lipson2021raftstereo}, shown in Figure
\ref{fig:raftstereo}, introduced a different paradigm: it adapted the
iterative refinement of RAFT optical flow to stereo. It builds an all-pairs
one-dimensional correlation volume, uses a separate context encoder to
seed a multi-level convolutional gated recurrent unit, and refines the
disparity through many lightweight update iterations, each of which looks
up correlation values around the current estimate. It restores full
resolution with a learned convex upsampling that predicts a weighted
combination of a coarse neighborhood.

\begin{figure}[!htbp]
\centering
\includegraphics[width=\textwidth]{fig_2_4b_raftstereo}
\caption[RAFT-Stereo architecture, representing the iterative refinement paradigm.]{RAFT-Stereo architecture, representing the iterative refinement paradigm \cite{lipson2021raftstereo}.}
\label{fig:raftstereo}
\end{figure}

Three of these mechanisms, the group-wise cost lookup, the recurrent
update, and convex upsampling, reappear in modified form in our own
design. IGEV-Stereo
\cite{xu2023igev} improved the paradigm by building a geometry encoding
volume, a lightweight three-dimensional-regularized cost volume that both
provides a warm-start disparity through soft-argmin and is looked up each
iteration alongside the all-pairs correlation, as in \eqref{eq:gev_lit};
with this warm start it matches RAFT-Stereo's thirty-two iterations in
three. The geometry encoding volume is the direct ancestor of the geometry
branch we use. The warm start matters because an iterative method spends its early
iterations recovering what a good initialization provides for free; with
it, IGEV-Stereo at a single iteration already outperforms RAFT-Stereo by
a wide margin, and the auxiliary loss on the initial estimate forces the
geometry volume to be a genuine disparity predictor rather than a mere
feature bank. Two later works refined the recurrent core in directions
relevant to this thesis. CREStereo \cite{li2022crestereo} introduced a
cascaded recurrent refinement whose adaptive group correlation layer
searches a small two-dimensional window rather than assuming a perfectly
horizontal epipolar line, and its ablation showed local correlation
outperforming the all-pairs form outright; both findings matter here,
the first because a two-dimensional local search is precisely what
confers tolerance to imperfect rectification, the second because it
justifies replacing the all-pairs volume with a small local lookup in a
compact model. Selective-Stereo \cite{wang2024selective} added a
selective recurrent unit that runs two gate branches with different
kernel sizes and fuses them per pixel through a contextual attention
map, reaching the quality of many more iterations in a few; when we tried
the same unit in our ablation campaign it did not transfer to the
compact tile setting.

The recurrent update is a convolutional gated recurrent unit
\cite{lipson2021raftstereo} that updates a hidden state from a
per-iteration input,
\begin{equation}
\begin{aligned}
z_n &= \sigma(\mathrm{Conv}([h_{n-1}, x_n])), \qquad
r_n = \sigma(\mathrm{Conv}([h_{n-1}, x_n])),\\
\tilde{h}_n &= \tanh(\mathrm{Conv}([r_n \odot h_{n-1}, x_n])), \qquad
h_n = (1 - z_n)\odot h_{n-1} + z_n \odot \tilde{h}_n
\end{aligned}
\label{eq:convgru_lit}
\end{equation}
where $x_n$ assembles the correlation lookup, the encoded disparity, and
the context, $h_{n-1}$ is the previous hidden state, $z_n$ and $r_n$ the
update and reset gates, $\tilde{h}_n$ the candidate state, $\sigma$ the
sigmoid, and $\odot$ the element-wise product; a small head then predicts
the disparity residual from $h_n$.

The geometry encoding volume of IGEV-Stereo \cite{xu2023igev} regularizes a group-wise correlation volume
and combines it with the all-pairs correlation into a lookup structure,
\begin{equation}
C_G = R\big(C_{\mathrm{gwc}}(f_{L}, f_{R})\big), \qquad
d_0 = \sum_{d=0}^{D-1} d \cdot \operatorname{softmax}_d\!\big(C_G(d)\big)
\label{eq:gev_lit}
\end{equation}
where $C_{\mathrm{gwc}}$ is the group-wise correlation of the
quarter-resolution features, $R$ a lightweight three-dimensional network,
$C_G$ the regularized geometry volume, and $d_0$ the warm-start disparity;
each iteration then looks up values from both $C_G$ and the all-pairs
correlation.

A parallel line replaced the cost volume itself with attention. STTR
\cite{li2021sttr} treats matching as dense sequence-to-sequence
correspondence along the epipolar lines, alternating self-attention
within each image and cross-attention between the views, with an
optimal-transport assignment that handles occlusion explicitly and
removes the fixed maximum-disparity assumption altogether. The
transformer family improved occlusion reasoning and freed the disparity
range, but its dense attention has remained costly at the resolutions
and budgets that edge deployment demands, which is why we draw on the
recurrent family instead.

\FloatBarrier

% The timeline is intentionally a full-page rotated figure so that its labels
% remain readable at print size.  Placing it at the section transition avoids
% leaving the lower half of the preceding page empty.
% Keep the established numbering: the foundation architecture remains Fig. 2.9
% and this timeline remains Fig. 2.10, even though the timeline is typeset first.
\setcounter{figure}{9}
\begin{sidewaysfigure}
\centering
\includegraphics[width=\textheight]{fig_2_2_timeline_preview}
\caption{Evolution of deep stereo matching, 2016 to 2026.}
\label{fig:timeline}
\end{sidewaysfigure}
\setcounter{figure}{8}

\subsection{Foundation-Model Era}

\begin{figure}[H]
\centering
\includegraphics[width=\textwidth]{fig_2_foundation_arch}
\caption[FoundationStereo architecture, representing the foundation-model-era paradigm.]{FoundationStereo architecture, representing the foundation-model-era paradigm \cite{wen2025foundationstereo}.}
\label{fig:foundation_arch}
\end{figure}
\setcounter{figure}{10}

The most recent and most accurate systems follow a similar approach: they
graft a frozen monocular vision foundation model onto an iterative stereo
core and train on very large data. FoundationStereo
\cite{wen2025foundationstereo} fuses a frozen DepthAnythingV2 transformer
with a light convolutional encoder through a side-tuning adapter, adds an
attentive hybrid cost filter, and trains on a million synthetic pairs,
reaching 0.34 pixel end-point error on SceneFlow; its architecture is
shown in Figure \ref{fig:foundation_arch}. DEFOM-Stereo
\cite{jiang2025defom} adds a monocular depth foundation model to a
RAFT-Stereo core and splits each iteration into a scale update that
corrects the monocular prior and a standard refinement. MonSter
\cite{cheng2025monster} couples an IGEV stereo branch with a monocular
branch through bidirectional refinement, and Stereo-Anywhere
\cite{bartolomei2025stereoanywhere} uses a monocular prior to truncate the
stereo cost volume and suppress the false matches that mirrors and glass
produce. The prior enters each method differently. DEFOM-Stereo splits
each recurrent iteration into a scale update, which corrects the
monocular prior's per-pixel scale through a dedicated lookup over
candidate scale factors, and a standard residual update; its ablation
shows the monocular initialization alone hurts and only becomes useful
with the scale correction. MonSter runs the two branches symmetrically,
each guiding the other after a global scale and shift alignment, and its
ablation shows four iterations with the prior matching what the plain
stereo core needs thirty-two iterations to reach. The accuracy gain
therefore comes from two sources: the prior's knowledge of monocular
structure and the iterations it saves. These systems set the accuracy
ceiling, but they do so with hundreds of millions of parameters and
hundreds of milliseconds of latency, a cost that motivates the edge
problem we address rather than solving it. Figure
\ref{fig:timeline} traces this progression across the six paradigm eras
and marks where the present work sits.

\FloatBarrier
\subsection{Requirements of Edge-Deployable Stereo Matching}
\label{sec:edge_requirements}
Deployment on an edge accelerator cannot be solved by choosing the most
accurate published network: models at the accuracy ceiling carry
hundreds of millions of parameters and hundreds of milliseconds of
latency, and mid-size iterative networks run their update loop tens of
times per frame on desktop graphics hardware. Edge accelerators on
mobile robots and drones, exemplified by the Jetson Orin Nano class of
device, impose four requirements that bound every design choice we make
in this thesis.

First, the network must be compact enough for on-device memory:
accelerators share memory with the robot's other software, so the
parameter budget must stay within a few million parameters, the region
occupied by the efficient family surveyed next. Second, the
network must run in real time on a device of roughly four to six
tera-operations per second; a perception stage that cannot keep up
with the camera forces the robot to act on stale depth, so latency is
a hard requirement that rules out many-iteration update loops and
dense attention. Third, the network must keep a small inference
working set. Peak activation memory, dominated by the cost volume in
most stereo networks, must fit alongside the model weights and the
application; this is why cost-volume compression forms its own
technique family below. Fourth, the network must use operators that
quantize cleanly to eight-bit integers, since accelerators reach rated
throughput only under integer quantization; operators that resist
quantization, such as trilinear interpolation in three-dimensional
volumes, forfeit that throughput regardless of parameter count. These
observations frame the classification below.

\FloatBarrier
\subsection{Classification of Efficient Stereo Methods}
\label{sec:classification}
The efficient and edge-deployable methods are the closest relatives of the
model proposed in this thesis, and they are best understood by the
technique family each uses to reduce cost. Figure \ref{fig:taxonomy}
organizes these families, and the subsections below treat each in turn.

Figure \ref{fig:reqmap} makes the connection explicit: each deployment
requirement is answered primarily by one or two of the technique
families, and a family that answers no requirement of the target
platform can be skipped without loss.

\begin{figure}[!htbp]
\centering
\includegraphics[width=0.95\textwidth]{fig_2_7_reqmap_preview}
\caption{Edge-deployment requirements mapped to compression technique families.}
\label{fig:reqmap}
\end{figure}

\begin{figure}[!t]
\centering
\includegraphics[width=0.82\textwidth]{fig_2_5_taxonomy_preview}
\caption{Taxonomy of compression techniques for deep stereo matching.}
\label{fig:taxonomy}
\end{figure}

\subsection{Backbone Substitution}
The simplest way to reduce cost is to replace the heavy feature encoder
with a mobile backbone. MobileStereoNet
\cite{shamsafar2022mobilestereonet} extended MobileNet's separable
blocks into the three-dimensional cost-volume regularization itself and
reported two design points: a two-dimensional variant at 1.14 pixel
end-point error with 2.23 million parameters, and a three-dimensional
variant at 0.80 pixels with 1.77 million, an order of magnitude smaller
than the contemporaneous volumetric baselines at comparable accuracy.
The lesson generalizes beyond MobileNet: any well-pretrained compact
backbone can serve, and a detection backbone carries an added advantage,
since its features are already trained to localize objects across
scales. We follow this route and use a truncated detection backbone as
our shared encoder.

\subsection{Cost-Volume Compression}
The three-dimensional cost volume is the dominant cost in most stereo
networks, and several families compress it. Cascade methods
\cite{gu2020cascade} build a coarse volume over the full disparity range
and then progressively narrow the range at finer scales; on the cascade
paper's reproduced PSMNet baseline the two-stage cascade cut memory
substantially while also improving the end-point error from 0.887 to
0.721 pixels (a reproduced baseline, distinct from the 1.09 pixels PSMNet
reports in Table \ref{tab:method_comparison}), so narrowing the range
turned out to be cheaper and more accurate at once. BGNet
\cite{xu2021bgnet} learns a bilateral grid that upsamples a
one-eighth-resolution cost volume in an edge-preserving way, with a
parameter-free slicing step, at 25 milliseconds per frame. DeepPruner
\cite{duggal2019deeppruner} predicts a per-pixel disparity range with a
differentiable PatchMatch and a confidence-range predictor, so the
expensive aggregation runs only on a pruned set of hypotheses; its fast
variant reaches 0.97 pixel end-point error at 61 milliseconds. HITNet
\cite{tankovich2021hitnet}, shown in Figure \ref{fig:hitnet}, removes the
three-dimensional cost volume altogether: it represents the scene as a
grid of small tiles, each carrying a slanted-plane hypothesis with a
disparity, two slopes, and a descriptor, and refines these hypotheses
through two-dimensional propagation and warping across scales, reaching
0.43 pixel end-point error at 0.97 million parameters and 54
milliseconds. The slanted plane lets a tile represent sub-pixel structure
that a fronto-parallel hypothesis cannot, and this is what lets the
network stay accurate without a dense volume. This tile-hypothesis
representation, including its plane slopes, is the paradigm from which we
build our model.

\begin{figure}[!htbp]
\centering
\includegraphics[width=\textwidth]{fig_2_4c_hitnet}
\caption[HITNet architecture, representing the tile-hypothesis paradigm.]{HITNet architecture, representing the tile-hypothesis paradigm \cite{tankovich2021hitnet}.}
\label{fig:hitnet}
\end{figure}

\subsection{Iterative-Loop Compression}
The update loop is where iterative methods spend most of their cost, and
the efficient variants target it directly. LightStereo
\cite{guo2025lightstereo} replaces three-dimensional aggregation with a
two-dimensional network of inverted-residual blocks that expand the
disparity channels, reaching 0.73 pixel end-point error at 3.44 million
parameters and 17 milliseconds; its authors observe that feature
extraction, not aggregation, dominates the remaining runtime.
Distill-then-Prune \cite{pan2024dtp} designs a backbone that is compatible
with deployment toolchains by construction: it avoids three-dimensional
convolution and trilinear interpolation, and reaches 1.56 pixel end-point
error at 0.26 million parameters with 16 milliseconds of latency on an
embedded board.

\subsection{Knowledge Distillation}
A distinct family reduces the accuracy cost of a small network by training
it against a large teacher. LiteAnyStereo \cite{jing2025liteanystereo},
whose network is shown in Figure \ref{fig:liteany}, pairs a 7.6 million
parameter non-iterative network with a three-stage training pipeline:
supervised training on synthetic data, self-distillation under input
perturbation, and distillation from a frozen FoundationStereo teacher on
unlabeled real pairs. The result is strong zero-shot accuracy that the
same network cannot reach by supervised training alone.

\begin{figure}[!htbp]
\centering
\includegraphics[width=\textwidth]{fig_2_4d_liteanystereo}
\caption[LiteAnyStereo architecture, representing the distillation-trained efficient paradigm.]{LiteAnyStereo architecture, representing the distillation-trained efficient paradigm \cite{jing2025liteanystereo}.}
\label{fig:liteany}
\end{figure}

Distill-then-Prune \cite{pan2024dtp} shows that distilling before pruning,
and using the teacher output alone rather than the teacher combined with
ground truth, gives the best student. Written generically, rather than
in any one paper's notation, the distillation term of such pipelines is
a masked regression between the student and teacher disparity,
\begin{equation}
\mathcal{L}_{\mathrm{KD}} = \frac{1}{|\Omega|}
  \sum_{i \in \Omega} m_i \cdot
  \operatorname{smooth}_{L_1}\!\big(\hat{d}_S(i) - \hat{d}_T(i)\big)
\label{eq:kd_loss}
\end{equation}
where $\hat{d}_S$ and $\hat{d}_T$ are the student and teacher disparities,
$\Omega$ the image domain, and $m_i$ an optional teacher-confidence mask.
This three-stage pipeline, illustrated in Figure \ref{fig:kd_stages}, has
not yet been applied to our model, although it may offer
the largest accuracy gain among the surveyed mechanisms and remains
promising future work.

\begin{figure}[!htbp]
\centering
\includegraphics[width=\textwidth]{fig_2_7_kd_stages}
\caption[The three-stage knowledge-distillation pipeline of LiteAnyStereo.]{The three-stage knowledge-distillation pipeline of LiteAnyStereo \cite{jing2025liteanystereo}.}
\label{fig:kd_stages}
\end{figure}

\subsection{Architectural Compression and Adaptive Compute}
CoEx \cite{bangunharcana2021coex} re-enriches a cheap,
information-poor correlation volume from the image side: its guided
cost-volume excitation derives per-channel attention weights from the
left features with cheap convolutions and applies them at every
aggregation scale, recovering the accuracy of a far costlier
concatenation volume at 2.72 million parameters and 27 milliseconds. It
also swaps the full soft-argmin for a regression over the top disparity
candidates, avoiding the between-modes averaging of
\eqref{eq:softargmin}. CGI-Stereo \cite{xu2023cgistereo} fuses context
into geometry at every decoder scale with bidirectional gradient flow,
beating far larger networks zero-shot at a fraction of their cost. The
same trade, cheap image-derived guidance for volumetric capacity,
recurs in our guided excitation and mask-based upsampling.
LightStereo \cite{guo2025lightstereo} adds a complementary finding:
expanding the disparity channels in its inverted-residual aggregation
blocks helps, but expanding the spatial kernels does not, so the useful
capacity lies along the disparity axis.

Adaptive-compute methods change the cost at run time. AnyNet
\cite{wang2019anynet} produces a usable disparity from a first coarse
stage and refines it through residual cost volumes with a narrow search
range, letting an embedded board trade 29 milliseconds of coarse depth
against 97 milliseconds of refined depth as the budget allows. MADNet
\cite{tonioni2019madnet} instead adapts online to a new environment,
updating one module per frame under a photometric self-supervised loss
via a bandit strategy for module choice, and recovers most of the
accuracy a zero-shot deployment loses while sustaining 25 frames per
second. AnyNet's narrow-range search and both methods' coarse-to-fine
budgeting anticipate the local lookup and iteration schedule we use.
Neural architecture search reaches the same pipeline in LEAStereo
\cite{cheng2020leastereo}, which jointly searches the feature and
matching networks within the stereo pipeline and reaches 0.78 pixel
end-point error at 1.81 million parameters, showing a well-allocated
small budget can rival hand designs an order of magnitude larger.

Three lessons from the efficient family shape our design.
First, the cost volume is where the budget goes:
every successful compact method shrinks, prunes, or replaces it with a
sparser hypothesis structure, and the tile representation we adopt is the
most aggressive. Second, image-derived guidance is cheap and effective:
excitation, bilateral-grid, and mask-based upsampling all steer a small
volume or coarse prediction from the image, a pattern we follow with our
convex upsampling and context stream. Third, the training recipe is
not an afterthought: the distillation results show the same small
architecture can land at very different accuracy purely from how it is
trained, so we treat and ablate our training protocol with the same
rigor as the architecture. These families are largely
complementary; our model combines several, pairing a substituted
backbone with a compressed cost volume and a compressed recurrent loop.

\FloatBarrier
\section{Research Gap}
\label{sec:research_gap}
The survey above points to a specific opening rather than a general call
for smaller networks. Three observations, taken together, define it.

First, accuracy and deployability sit at opposite ends of the surveyed
work. The foundation-model and iterative systems hold the accuracy
ceiling but carry tens to hundreds of millions of parameters and hundreds
of milliseconds of latency, which no edge accelerator of the target class
can afford. The efficient family fits the budget, but each of its members
attacks a single source of cost in isolation: backbone substitution
replaces the encoder, cost-volume compression shrinks or removes the
volume, iterative-loop compression trims the update schedule, and
knowledge distillation reduces the accuracy penalty through training
rather than architecture. What is missing is a single compact design that
combines several of these reductions at once, a substituted detection
backbone, a tile-hypothesis representation in place of a dense volume, a
short recurrent loop, and a narrow-band geometry volume, and shows that
the combination still lands inside the edge budget. Building and
justifying that combination is the task of Chapter~\ref{ch:design}.

Second, the efficient literature is reported almost entirely on
in-domain accuracy, parameters, and latency, and says far less about how
these compact models behave when the input geometry is imperfect. The
benchmarks that drive the field are near-perfectly rectified, yet the
low-cost stereo rigs these models are meant for leave a residual vertical
misalignment of up to about a pixel. A method that is only ever measured
on flawless input leaves its field robustness untested, and the survey
found no compact model that reports its accuracy as a function of
controlled rectification error. Closing that measurement gap is the
second objective of this thesis.

Third, the survey left one methodological question unsettled: how much of
a compact network's cross-domain accuracy comes from its architecture and
how much from its training recipe. The distillation results show the same
small network reaching very different zero-shot accuracy purely from how
it is trained, while other efficient methods credit their recurrent core.
The literature does not resolve this with controlled evidence, which is
why this thesis treats the training protocol as a first-class design
variable and ablates it under the same single-knob discipline it applies
to the architecture.

These three gaps, an unbuilt combination, an untested robustness, and an
unsettled attribution, are what the comparison in the next section
quantifies and what the design chapter sets out to close.

\FloatBarrier
\section{Comparative Analysis}
\label{sec:comparative_analysis}

\subsection{Cross-Domain Generalization}
A network trained on synthetic data often fails on real domains: it
learns shortcut features tied to the synthetic statistics, such as
consistent local color distributions, instead of domain-invariant
structure, and it breaks down in textureless, reflective, repetitive, and
large-disparity regions. The scale of the failure is substantial: GA-Net
trained on synthetic data and evaluated zero-shot reaches 11.7 percent
outliers on KITTI 2015 and 32.2 percent on Middlebury, and PSMNet reaches
16.3 percent on KITTI, several times their fine-tuned rates
\cite{zhang2020dsmnet}. Two families
of remedy exist. Architecture-level methods make the features themselves
domain-invariant. DSMNet \cite{zhang2020dsmnet} adds a domain
normalization that removes image-level and pixel-level statistic shifts
together with a structure-preserving graph filter, and improves the
zero-shot KITTI rate from 11.7 to 6.5 percent and the Middlebury rate
from 32.2 to 21.8 percent. HVT \cite{chang2023hvt} applies hierarchical
visual transformations at global, local, and pixel level during training
only, at zero inference cost, and pushes the zero-shot KITTI rate further
to 4.9 percent on the same backbone. Training-recipe methods change the
data and the supervision instead: StereoAnything
\cite{guo2024stereoanything} mixes labeled stereo with tens of millions
of monocular-to-stereo pseudo-pairs and shows that dataset diversity
matters more than dataset size, and LiteAnyStereo
\cite{jing2025liteanystereo} reaches strong zero-shot accuracy with a
non-iterative network trained by three-stage distillation. The relative
importance of iteration and training pipeline for generalization is not
settled in the literature: some efficient methods argue that iteration is
necessary for cross-domain accuracy, while LiteAnyStereo reaches
competitive zero-shot numbers without any iteration, and the cleanest
controlled evidence attributes its accuracy to the training pipeline
rather than to any recurrent module. Our own non-iterative tile network
follows this second line of evidence.

\FloatBarrier
\subsection{Benchmarks and Datasets}
Table \ref{tab:datasets} summarizes the benchmarks used in this thesis.
SceneFlow is a large synthetic dataset of 35,454 training pairs across
three subsets with dense rendered ground truth, and it is the standard
pretraining set; only its FlyingThings3D subset provides an official
test split, which is used here for in-domain evaluation. Its value lies
in the density and exactness of its labels, which real sensors cannot
provide, at the price of a domain gap discussed in the preceding
subsection. KITTI 2012 and 2015 provide real driving scenes with sparse
LiDAR ground truth, and their D1-all outlier rate has become the
field's standard metric for driving. ETH3D covers mixed indoor and
outdoor scenes at gray-scale with laser-scanned ground truth. Middlebury
2014 provides high-resolution indoor scenes with dense structured-light
ground truth, and it is the hardest of the four in two respects: its
images reach several megapixels, so most methods evaluate at half or
quarter resolution, and its disparities extend to several hundred pixels,
far beyond the search range most networks are built for. These two
properties make it the standard zero-shot stress test, and it is the
cross-domain benchmark we use. Its evaluation protocol splits the
scenes into a training set with public ground truth and a test set whose
ground truth is withheld on a server. The zero-shot convention we follow,
matching the reference methods we compare against, evaluates on the
ground-truth-public scenes without any training or fine-tuning on
Middlebury. The metrics also carry conventions worth
noting: end-point error rewards sub-pixel accuracy on smooth surfaces and
can look healthy while boundaries are poor, whereas the bad-$t$ and
D1-all rates count outlier pixels and expose exactly the boundary and
textureless failures that end-point error averages away. For this
reason, we report both families of metric together.

% Place the two-part Pareto figure before the comparison table so the open
% space on the page is used by the figures rather than left below the table.
\begin{figure}[H]
\centering
\includegraphics[width=0.94\textwidth]{fig_2_6a_param_pareto}
\caption{Accuracy against parameters and latency for the surveyed methods: (a) accuracy against trainable parameters.}
\label{fig:pareto}
\end{figure}

\begin{figure}[H]
\ContinuedFloat
\centering
\includegraphics[width=0.94\textwidth]{fig_2_6b_latency_pareto}
\caption[]{Accuracy against parameters and latency for the surveyed methods (continued): (b) KITTI 2015 D1-all against reported latency.}
\end{figure}

\begin{table}[!htbp]
\centering
\caption{Stereo benchmark datasets with train/test splits, resolutions, and ground-truth types.}
\label{tab:datasets}
\small
\begin{tabular}{lllllc}
\toprule
Dataset & Domain & Train/Test & Resolution & Ground truth & Ref.\\
\midrule
SceneFlow & synthetic & 35{,}454 / 4{,}370 & $960{\times}540$ & dense, rendered & \cite{mayer2016sceneflow}\\
KITTI 2012 & driving & 194 / 195 & $1242{\times}375$ & sparse LiDAR & \cite{geiger2012kitti}\\
KITTI 2015 & driving & 200 / 200 & $1242{\times}375$ & sparse LiDAR + CAD & \cite{menze2015kitti}\\
Middlebury 2014 & indoor & 23 / 10 & up to ${\sim}3000{\times}2000$ & structured light & \cite{scharstein2014middlebury}\\
ETH3D & mixed & 27 / 20 & ${\sim}940{\times}490$ & laser scan & \cite{schops2017eth3d}\\
\bottomrule
\end{tabular}
\end{table}

\FloatBarrier
\subsection{Accuracy, Parameters, and Latency Across Method Families}
Table \ref{tab:method_comparison} and Figure \ref{fig:pareto} place
every method surveyed in this chapter, plus our own model, on a common
footing of parameters, SceneFlow end-point error, KITTI 2015 D1-all, and
reported latency. Read together with the taxonomy of
Section~\ref{sec:classification}, they turn the family-by-family survey
above into a direct comparison. Foundation-scale and iterative systems
hold the accuracy frontier, at hundreds of millions and tens of millions
of parameters respectively, while the efficient family occupies the
low-parameter region at some cost in accuracy: FoundationStereo reaches
0.34 pixel end-point error at roughly 340 million parameters and 470
milliseconds, whereas HITNet reaches 0.43 pixels, a difference of only
0.09 pixels, at 0.97 million parameters and 54 milliseconds, under a
hundredth of the latency. Parameter count alone does not order the
methods either: PSMNet, at 5.2 million parameters, trails HITNet's 0.97
million on both accuracy and latency, which shows that the architecture
a budget buys matters as much as the size of the budget. CoEx and
LightStereo-S sit in the same low-parameter, low-latency region as
HITNet, at 2.72 and 3.44 million parameters and 27 and 17 milliseconds,
and land within 0.3 pixels of it, so several distinct efficient
architectures reach comparable operating points from different
technique families. Our own model, StereoLite, sits at 2.96 million
parameters and 49.8 milliseconds with 0.78 pixel end-point error,
inside the same efficient cluster but with room to close against HITNet
on accuracy.

The accuracy-parameter-latency comparison alone is not sufficient,
however, because it says nothing about cross-domain behavior. The review
in Section~\ref{sec:existing_research} showed that two networks of
similar synthetic-data accuracy can differ by an order of magnitude in
zero-shot outlier rate depending on their training pipeline rather than
their parameter count, so a fair comparison across families has to weigh
generalization alongside the raw numbers in Table
\ref{tab:method_comparison}. Our own design targets this gap directly,
sizing the architecture against the frontier established here.

\begin{table}[!htbp]
\centering
\caption{Parameters, SceneFlow error, KITTI D1-all, and latency for representative methods (values verified against cited papers).}
\label{tab:method_comparison}
\small
\begin{tabular}{lccccc}
\toprule
Method & Year & Params (M) & SF EPE (px) & KITTI D1 (\%) & Latency (ms)\\
\midrule
PSMNet \cite{chang2018psmnet} & 2018 & 5.2 & 1.09 & 2.32 & 410\\
HITNet (L) \cite{tankovich2021hitnet} & 2021 & 0.97 & 0.43 & 1.98 & 54\\
BGNet \cite{xu2021bgnet} & 2021 & ${\sim}2.9$ & 1.17 & 2.51 & 25\\
CoEx \cite{bangunharcana2021coex} & 2021 & 2.72 & 0.69 & 2.13 & 27\\
RAFT-Stereo \cite{lipson2021raftstereo} & 2021 & 11.23 & 0.61 & 1.82 & 380\\
IGEV-Stereo \cite{xu2023igev} & 2023 & 12.60 & 0.47 & 1.59 & 180\\
LightStereo-S \cite{guo2025lightstereo} & 2025 & 3.44 & 0.73 & 2.30 & 17\\
DEFOM-Stereo \cite{jiang2025defom} & 2025 & 47.3 & 0.42 & 1.55 & 316\\
FoundationStereo \cite{wen2025foundationstereo} & 2025 & ${\sim}340$ & 0.34 & 1.46 & ${\sim}470$\\
\midrule
StereoLite (ours) & 2026 & 2.96 & 0.78 & n/a & 49.8\\
\bottomrule
\end{tabular}
\end{table}

\FloatBarrier
\section{Summary}
\label{sec:ch2_summary}
This chapter surveyed the literature that bears on an edge-deployable
stereo network and then compared that literature directly rather than
treating it as a list of isolated methods. We traced deep stereo
matching from the classical four-step pipeline through learned matching
costs, three-dimensional cost-volume networks, iterative refinement, and
the foundation-model systems that now hold the accuracy ceiling at a
computational cost most edge devices cannot meet. From that survey we
stated four requirements for edge deployment (compact parameters,
real-time latency on roughly four to six tera-operations per second, a
small inference working set, and integer-friendly operators) and used
them to classify the efficient literature into backbone substitution,
cost-volume compression, iterative-loop compression, knowledge
distillation, and architectural and adaptive-compute techniques. We
compared these families directly: the review of cross-domain
generalization showed that both architecture-level and training-recipe
remedies narrow the synthetic-to-real gap, and that the relative
contribution of iteration versus training pipeline to zero-shot accuracy
remains unsettled in the literature; the benchmark review fixed the
dual-metric, zero-shot evaluation convention the rest of this thesis
follows; and the combined parameter-accuracy-latency comparison placed
our own model on the same footing as every method it is compared
against.

% SOURCES: thesis/book/mainmatter/ch2.tex sections 2.1 (Introduction),
%   2.2 (Literature Review, all 5 subsections: Stereo Matching
%   Fundamentals, The Deep Learning Transition, The 3D Cost-Volume Era,
%   Iterative Refinement, Foundation-Model Era), 2.3 (Requirements of
%   Edge-Deployable Stereo Matching), 2.4 (Classification of Efficient
%   Stereo Methods, all 5 subsections: Backbone Substitution,
%   Cost-Volume Compression, Iterative-Loop Compression, Knowledge
%   Distillation, Architectural Compression and Adaptive Compute), 2.5
%   (Cross-Domain Generalization and Benchmarks, both subsections), and
%   the comparison Table \ref{tab:method_comparison} + Figure
%   \ref{fig:pareto} + their framing paragraph lifted out of old 2.6.
%   All \label names, equations, figures, and tables carried over
%   verbatim (byte-identical numbers/citations). Old top-level sections
%   2.3 and 2.4 were flattened into \subsection entries directly under
%   the new 2.2 "Existing Research" (no sub-subsections, per the
%   x.y.z.w ban), and old 2.5's two subsections became \subsection
%   entries under the new 2.3 "Comparative Analysis" alongside a new
%   third subsection ("Accuracy, Parameters, and Latency Across Method
%   Families") that reframes the comparison table/figure with fresh,
%   non-mechanism-specific comparative prose.
% DROPPED: nothing substantive; every paragraph, figure, table, and
%   equation from old 2.1-2.5 was carried over. Prose was lightly
%   rewritten in "we" voice where it referred to our own design choices,
%   and every hard-coded "Chapter N" was replaced with a \ref to the new
%   chapter label (\ref{ch:design} for the new Ch3, \ref{ch:results} for
%   the new Ch6 "Results and Discussion", \ref{ch:conclusion} for the
%   new Ch9). These three labels are assumed names for the sibling
%   chapter agents' \label{} declarations; if a sibling agent chose a
%   different label string, the QA pass (B4) will surface it as an
%   undefined reference and it should be reconciled at integration.
% MOVED-AWAY: old ch2 Section 2.6 "Detailed Review of the Adopted
%   Methodology" is owned by the ch3 agent (Engineering Design and
%   Methodology) per RESTRUCTURE_PLAN_V2. Specifically excluded from
%   this file: the four paragraphs describing the group-wise
%   correlation cost, ConvGRU + convex upsampling, the geometry encoding
%   volume, and the slanted tile-plane representation as "adopted
%   mechanisms" of the thesis's own network, and the closing paragraph
%   naming the specific combination as an unaddressed gap in the
%   surveyed literature. Only the comparison Table \ref{tab:method_comparison},
%   Figure \ref{fig:pareto}, and their general (non-mechanism-specific)
%   framing paragraph were retained here, per the explicit instruction
%   that "capability-matrix / comparison-table material" belongs in the
%   new 2.3.
